\input{sol_pre.tex}
\input{q_2026-03-07_1.tex}
\begin{document}
\input{sol_head.tex}

%================= 1
\Savol{1}
$24^3=2^9\cdot3^3$, \ $6^{12}=2^{12}\cdot3^{12}$, \ $20^3=2^6\cdot5^3$ $\Rightarrow$ $m=2^{27}\cdot3^{15}\cdot5^3$, \ $m^2=2^{54}\cdot3^{30}\cdot5^6$.\\
$\dfrac{m^2}{2^n}$ natural son bölişi uçun $n\le54$.
\Javob{A) 54}

%================= 2
\Savol{2}
Baholaş: $3<\sqrt{10}\Rightarrow a=\lg3<\frac12$; \ $\frac{5\pi}6<3<\pi\Rightarrow b=\cos3<-\frac{\sqrt3}2$; \ $c=\ln3>1$.\\
Demak $b<a<c$ va $a+b<0$, $c+a>0$:
\[\frac{(a-b)+(c-b)-(c-a)}{-(a+b)+(c-b)-(c+a)}=\frac{2a-2b}{-2a-2b}=\frac{b-a}{b+a}.\]
\Javob{B) $\frac{b-a}{b+a}$}

%================= 3
\Savol{3}
Oxiridan boşlaymiz. 3-ötişdan keyin $0$ qoldi: tölovdan oldin $36\,000$, ikki baravar bölişidan oldin $18\,000$.\\
2-ötiş: $(18\,000+36\,000):2=27\,000$. \quad 1-ötiş: $(27\,000+36\,000):2=31\,500$.
\Javob{D) 31\,500}

%================= 4
\Savol{4}
Neft ajratilgandan keyin $60-15=45$ litr qoldi, undagi neft $0{,}04\cdot45=1{,}8$ litr.\\
Boşida neft $1{,}8+15=16{,}8$ litr: \ $\dfrac{16{,}8}{60}\cdot100\%=28\%$.
\Javob{A) 28\%}

%================= 5
\Savol{5}
\[\frac{a^{-6}\cdot\left(-a^{14}\right)}{a^6}=-a^{-6+14-6}=-a^2.\]
\Javob{A) $-a^2$}

%================= 6
\Savol{6}
\[\frac1{2+\sqrt2}-\frac1{2-\sqrt2}=\frac{(2-\sqrt2)-(2+\sqrt2)}{4-2}=-\sqrt2,\qquad -\sqrt2-\sqrt2=-2\sqrt2.\]
\Javob{B) $-2\sqrt2$}

%================= 7
\Savol{7}
$P=\sqrt3+\sqrt{15}$, \ $Q=\sqrt7+\sqrt{11}$: maxrajlar $P+Q$ va $P-Q$.
\[8(\sqrt3-\sqrt{15})\cdot\frac{2P}{P^2-Q^2}=\frac{16(3-15)}{(18+6\sqrt5)-(18+2\sqrt{77})}=\frac{-192}{6\sqrt5-2\sqrt{77}}=\frac{-96}{3\sqrt5-\sqrt{77}}.\]
\[\frac{-96\left(3\sqrt5+\sqrt{77}\right)}{45-77}=3\left(3\sqrt5+\sqrt{77}\right)=9\sqrt5+3\sqrt{77}.\]
\Javob{D) $9\sqrt5+3\sqrt{77}$}

%================= 8
\Savol{8}
Logarifmlar arifmetik progressiya $\Rightarrow$ $162,\ x,\ y,\ z,\ 1250$ — geometrik progressiya:
\[q^4=\frac{1250}{162}=\frac{625}{81}\ \Rightarrow\ q=\frac53\ \ (\text{hadlar musbat}).\]
$x=270$, \ $y=450$, \ $z=750$: \ $z-x-y=750-720=30$.
\Javob{C) 30}

%================= 9
\Savol{9}
$b_2^3=216\Rightarrow b_2=6$; \ $b_1+b_3=15$, \ $b_1b_3=36$ $\Rightarrow$ $\{3;12\}$. \ $q<1$: \ $12;\ 6;\ 3$, \ $q=\frac12$.
\[S_5=\frac{12\left(1-\frac1{32}\right)}{1-\frac12}=24\cdot\frac{31}{32}=\frac{93}4=23\frac14.\]
\Javob{A) $23\frac14$}

%================= 10
\Savol{10}
$(2a+b)^2-3a(a+b)=a^2+ab+b^2$, \ $a^3-b^3=(a-b)(a^2+ab+b^2)$ $\Rightarrow$ birinçi kasr $a-b$.\\
$(a^2+b)(b^2+a)-ab(ab+1)=a^3+b^3$ $\Rightarrow$ ikkinçi kasr $a+b$. \quad $(a-b)+(a+b)=2a$.
\Javob{B) $2a$}

%================= 11
\Savol{11}
$s=\sqrt{a^2-3a}$, \ $a^2-s^2=3a$:
\[\left(\frac{a+s}{a-s}-\frac{a-s}{a+s}\right)\cdot\frac34=\frac{(a+s)^2-(a-s)^2}{a^2-s^2}\cdot\frac34=\frac{4as}{3a}\cdot\frac34=s=\sqrt{a^2-3a}.\]
\Javob{A) $\sqrt{a^2-3a}$}

%================= 12
\Savol{12}
$b=\sin167^\circ=\sin13^\circ$, \ $c=\sin375^\circ=\sin15^\circ$, \ $a=\sin17^\circ$.\\
$[0^\circ;90^\circ]$ da sinus ösuvçi: $\sin17^\circ>\sin15^\circ>\sin13^\circ$, ya'ni $a>c>b$.
\Javob{A) $a>c>b$}

%================= 13
\Savol{13}
Hammasini bir tomonga ötkazib, töla kvadratlarga ajratamiz:
\[\left(2\tg^2x-4\tg x+2\right)+\left(2\sin^2x-2\sqrt2\sin x+1\right)=0\ \Rightarrow\ 2(\tg x-1)^2+\left(\sqrt2\sin x-1\right)^2=0.\]
$\tg x=1$ va $\sin x=\frac{\sqrt2}2$ $\Rightarrow$ $x=\frac\pi4+2\pi k$. \ $[0;4\pi]$ da: $\frac\pi4$, $\frac{9\pi}4$; \ yiğindi $\frac{10\pi}4=\frac{5\pi}2$.
\Javob{A) $\frac{5\pi}2$}

%================= 14
\Savol{14}
$\sqrt{7\pm4\sqrt3}=2\pm\sqrt3$, \ $(2+\sqrt3)(2-\sqrt3)=1$. \ $t=(2+\sqrt3)^x$:
\[t+\frac1t=14\ \Rightarrow\ t=7\pm4\sqrt3=(2\pm\sqrt3)^2\ \Rightarrow\ x=2\ \text{yoki}\ x=-2.\]
Yiğindi: $0$.
\Javob{D) 0}

%================= 15
\Savol{15}
Şart: $x>1$. \ $\left(\frac49\right)^{\log_{\frac19}(x-1)}=\left(\frac23\right)^{2\cdot\left(-\frac12\log_3(x-1)\right)}=\left(\frac23\right)^{-\log_3(x-1)}$:
\[\log_3(x+1)-\log_3(x-1)=1\ \Rightarrow\ \frac{x+1}{x-1}=3\ \Rightarrow\ x=2.\]
\Javob{A) 2}

%================= 16
\Savol{16}
$s=x+y$, \ $p=xy$: \ $s^2-2p=34$, \ $s+p=7$ $\Rightarrow$ $s^2+2s-48=0\Rightarrow s=6$ yoki $s=-8$.\\
$s=6,\ p=1$: $t^2-6t+1=0$ ($D>0$) — 2 ta yeçim, $x$ lar yiğindisi $6$.\\
$s=-8,\ p=15$: $t^2+8t+15=0\Rightarrow\{-3;-5\}$ — 2 ta yeçim, $x$ lar yiğindisi $-8$.
\[x_1+\dots+x_4=6+(-8)=-2.\]
\Javob{C) $-2$}

%================= 17
\Savol{17}
Aniqlaniş sohasi: $10+3x-x^2\ge0\Rightarrow -2\le x\le5$.\\
$\sqrt{\ }=0$: $x=-2$, $x=5$. \ $x^2-8x+12=0$: $x=2$, $x=6$ ($6$ sohaga kirmaydi).\\
Ildizlar: $-2;\ 2;\ 5$, \ yiğindisi $5$.
\Javob{A) 5}

%================= 18
\Savol{18}
1) $\dfrac{36-x^2}x\le0\Rightarrow\dfrac{(x-6)(x+6)}x\ge0\Rightarrow x\in[-6;0)\cup[6;+\infty)$.\\
2) $x^2+5x+25>0$, $|12-2x|\ge0$. Kasr $\le0$: yoki $|12-2x|=0$ ($x=6$), yoki $|x|-3<0$ ($-3<x<3$).\\
Kesişma: $(-3;0)\cup\{6\}$.
\Javob{C) $(-3;0)\cup\{6\}$}

%================= 19
\Savol{19}
Aniqlaniş sohasi $[11;12]$. Ikkala ildiz bir vaqtda nolga teng bölmaydi, şuning uçun yiğindi sohaning har bir nuqtasida musbat.
\Javob{B) $[11;12]$}

%================= 20
\Savol{20}
$f(-5)=8-5=3$, \ $f(3)=2-3=-1$, \ $f(-1)=8-1=7$.
\Javob{B) 7}

%================= 21
\Savol{21}
$x^2+x-6=(x+3)(x-2)$, \ $\dfrac{x+1}{(x+3)(x-2)}=\dfrac{2/5}{x+3}+\dfrac{3/5}{x-2}$:
\[\int_3^6=\frac25\ln\frac96+\frac35\ln\frac41=\frac25\ln\frac32+\frac65\ln2=\frac25\ln3+\frac45\ln2=\frac25\ln12.\]
\Javob{D) $\frac25\ln12$}

%================= 22
\Savol{22}
$S_n(x)=x+\dots+x^n$: \ $S_n(1)=n$, \ $S_n'(1)=\dfrac{n(n+1)}2$.
\[f'(1)=\frac{S_{2024}'(1)S_{1012}(1)-S_{2024}(1)S_{1012}'(1)}{S_{1012}(1)^2}=\frac{\frac{2024\cdot1012}2\,(2025-1013)}{1012^2}=\frac{2024}2=1012.\]
\Javob{D) 1012}

%================= 23
\Savol{23}
\fig{\begin{tikzpicture}[scale=0.32]
\coordinate (O1) at (0,5); \coordinate (O2) at (10.583,2);
\draw[thick] (-6,0)--(13.5,0);
\draw[thick] (O1) circle (5); \draw[thick] (O2) circle (2);
\coordinate (A) at ($(O1)+(58.3:5)$); \coordinate (B) at ($(O2)+(58.3:2)$); \coordinate (K) at ($(O1)+(58.3:3)$);
\draw[thick] ($(A)!-0.15!(B)$)--($(B)!-0.25!(A)$); \draw[thick,dashed] (O1)--(O2);
\draw[axs] (O1)--(A); \draw[axs] (O2)--(B); \draw[axs] (O2)--(K);
\fill (O1) circle (3pt) node[left]{\scriptsize$O_1$}; \fill (O2) circle (3pt) node[below]{\scriptsize$O_2$};
\node[above] at (A){\small$A$}; \node[above right] at (B){\small$B$}; \node[aux,left] at (K){\scriptsize$K$};
\end{tikzpicture}}
$O_1A\perp AB$, \ $O_2B\perp AB$. $O_2$ dan $AB$ ga parallel $O_2K$ ötkazamiz ($K\in O_1A$): $ABO_2K$ — töğri törtburçak, $O_1K=5-2=3$.
\[AB=O_2K=\sqrt{O_1O_2^2-O_1K^2}=\sqrt{121-9}=\sqrt{112}=4\sqrt7.\]
\Javob{A) $4\sqrt7$}

%================= 24
\Savol{24}
$x_1+x_2=2a$, \ $x_1x_2=-(a+6)$:
\[(x_1-x_2)^2=4a^2+4(a+6)=100\ \Rightarrow\ a^2+a-19=0\quad(D>0).\]
Viyet: $a_1+a_2=-1$.
\Javob{A) $-1$}

%================= 25
\Savol{25}
$\angle EAD=\angle EDA=t$ ($AE=DE$). $\angle DEC$ — $\triangle AED$ ning taşqi burçagi: $\angle DEC=2t$.\\
$DE=DC\Rightarrow \angle DCE=2t$; \ $AD=AC\Rightarrow\angle ADC=\angle ACD=2t$.\\
$\triangle ADC$: $t+2t+2t=180^\circ\Rightarrow t=36^\circ$, \ $\angle C=72^\circ$.\\
$AD$ — bissektrisa: $\angle A=2t=72^\circ$, \ $\angle B=180^\circ-72^\circ-72^\circ=36^\circ$.\\
$DF\perp AB\Rightarrow\angle BFD=90^\circ$, \ $FL$ — bissektrisa: $\angle BFL=45^\circ$.
\[\beta=\angle FLB=180^\circ-36^\circ-45^\circ=99^\circ.\]
\Javob{B) $99^\circ$}

%================= 26
\Savol{26}
$\triangle ADC$: $AD^2=8^2+5^2-2\cdot8\cdot5\cos60^\circ=49\Rightarrow AD=7$.\\
$\triangle ADC$ (8, 5, 7): $S=\frac12\cdot8\cdot5\cdot\frac{\sqrt3}2=10\sqrt3$, \ $p=10$, \ $r_1=\sqrt3$.\\
$\triangle ABD$ (8, 3, 7): $S=\frac12\cdot8\cdot3\cdot\frac{\sqrt3}2=6\sqrt3$, \ $p=9$, \ $r_2=\frac{2\sqrt3}3$.
\[\pi\left(r_1^2+r_2^2\right)=\pi\left(3+\frac43\right)=\frac{13\pi}3.\]
\Javob{C) $\frac{13\pi}3$}

%================= 27
\Savol{27}
$\triangle BOC$ va $\triangle COD$ ning $C$ dan balandligi umumiy: $\dfrac{BO}{OD}=\dfrac13$, $S_{COD}=9$.\\
Trapetsiyada $S_{AOB}=S_{COD}$ ($S_{ABD}=S_{ACD}$, umumiy qism $S_{AOD}$ ayriladi): $S_{AOB}=9$.
\Javob{C) 9}

%================= 28
\Savol{28}
Muntazam oltiburçakning qöşni tomonlari örtalari orasidagi masofa $a\cdot\frac{\sqrt3}2$ (120° li teng yonli uçburçak asosi). Yuzlar har safar $\frac34$ marta kamayadi.\\
$S_0=\frac{3\sqrt3}2$ dan boşlanuvçi çeksiz kamayuvçi geometrik progressiya yiğindisi:
\[S=\frac{S_0}{1-\frac34}=4\cdot\frac{3\sqrt3}2=6\sqrt3.\]
\Javob{C) $6\sqrt3$}

%================= 29
\Savol{29}
$\overrightarrow{AB}=(x-6;\ y-4;\ 2)\parallel(4;-3;1)$ $\Rightarrow$ $k=2$: \ $\overrightarrow{AB}=(8;-6;2)$, \ yiğindi $4$.
\Javob{A) 4}

%================= 30
\Savol{30}
Romb tomoni $\sqrt{30^2+40^2}=50$. Markaz $O$ dan tomongaça masofa (töğri burçakli uçburçak balandligi): $OK=\dfrac{30\cdot40}{50}=24$.\\
Uç perpendikulyar haqidagi teoremaga köra $PK\perp$ tomon:
\[PK=\sqrt{PO^2+OK^2}=\sqrt{45^2+24^2}=\sqrt{2601}=51.\]
\Javob{A) 51}

%================= 31
\Savol{31}
Barça qism töplamlar $2^6=64$; \ 3 ham, 5 ham qatnaşmaydiganlari $2^4=16$. \ $64-16=48$.
\Javob{B) 48}

%================= 32
\Savol{32}
$C_n^4-C_n^3=C_n^3\cdot\dfrac{n-7}4=14$. \ $n=8$: $56\cdot\frac14=14$ — mos ($n$ oşgani sari ifoda oşadi).\\
$C_8^2=28$.
\Javob{D) 28}

%================= 33-35
\newpage
\Savol{33}
\fig{\begin{tikzpicture}[scale=0.32]
\draw[thick] (0,0) circle (7.5); \fill (0,0) circle (3pt) node[below left]{\scriptsize$O$};
\draw[thick] (-3,6.874)--(3,6.874); \draw[axs] (0,0)--(0,6.874); \draw[ax] (0,0)--(3,6.874);
\node[aux,right] at (0,3.3){\scriptsize$d$}; \node[above] at (0,6.9){\scriptsize$6$};
\node[aux] at (2.5,3){\scriptsize$r$};
\node[below] at (0,-7.6){\small Asosga proyeksiya};
\end{tikzpicture}}
$2\pi r^2+2\pi r\cdot8=232{,}5\pi\Rightarrow r^2+8r-116{,}25=0\Rightarrow r=\dfrac{-8+23}2=7{,}5$.

\textbf{33.} Kesma uçlari ikki asos aylanalarida. Uning asosga proyeksiyasi — vatar, uzunligi $\sqrt{10^2-8^2}=6$. Kesma öqqa parallel tekislikda yotadi, öqdan masofa — markazdan vatargaça masofa:
\[d=\sqrt{7{,}5^2-3^2}=\sqrt{47{,}25}=\frac{3\sqrt{21}}2.\]
\textbf{34.} Kvadrat kesim tomoni $h=8$, ya'ni asosdagi vatar 8: \ $d=\sqrt{7{,}5^2-4^2}=\sqrt{40{,}25}=\dfrac{\sqrt{161}}2$.

\textbf{35.} $R=\sqrt{r^2+\left(\frac h2\right)^2}=\sqrt{56{,}25+16}=\sqrt{72{,}25}=8{,}5=\dfrac{17}2$.
\Javob{33 – D; \ 34 – A; \ 35 – C}

%================= 36
\newpage
\Savol{36}
Töla kvadratlarga ajratamiz:
\[(x-a)^2+(y+a)^2=(a-3)^2,\qquad (x+3a)^2+(y+4a)^2=(a-1)^2.\]
Markazlar $O_1(a;-a)$, $O_2(-3a;-4a)$, radiuslar $|a-3|$, $|a-1|$; markazlar orasi $d=\sqrt{16a^2+9a^2}=5|a|$.\\
$a=3$ yoki $a=1$ da bittasi nuqtaga aylanadi, lekin u ikkinçi aylanada yotmaydi ($15\ne2$, $5\ne2$); $a=0$ da konsentrik — yeçim yöq.\\
Yagona yeçim — aylanalar urinsa.\\
\textbf{Taşqi uriniş} $5|a|=|a-3|+|a-1|$: \ $a\in[0;1]$: $5a=4-2a\Rightarrow a=\frac47$; \ $a<0$: $-5a=4-2a\Rightarrow a=-\frac43$; \ $a>1$ da yeçim yöq.\\
\textbf{Içki uriniş} $5|a|=\bigl||a-3|-|a-1|\bigr|$: \ $a\le1$ da öng tomon $2$: $a=\pm\frac25$; \ $a>1$ da yeçim yöq.

a) $a\in\left\{-\frac43;\ -\frac25;\ \frac25;\ \frac47\right\}$ — 4 ta.

b) Eng kattasi $\frac47$.
\Javob{a) 4; \quad b) $\frac47$}

%================= 37
\newpage
\Savol{37}
$c=\cos2x$: \ $\sin^2x=\frac{1-c}2$, \ $\cos^2x=\frac{1+c}2$.
\[\sin^{10}x+\cos^{10}x=\frac{(1-c)^5+(1+c)^5}{32}=\frac{1+10c^2+5c^4}{16}.\]
\[1+10c^2+5c^4=29c^4\ \Rightarrow\ 24c^4-10c^2-1=0\ \Rightarrow\ c^2=\frac{10+14}{48}=\frac12\ \ \left(c^2=-\tfrac1{12}\ \text{emas}\right).\]
$\cos^22x=\frac12\Rightarrow\cos4x=0\Rightarrow x=\dfrac\pi8+\dfrac{\pi k}4$, $k\in\mathbb Z$.

a) Eng kiçik musbat ildiz: $\dfrac\pi8$.

b) Ildizlar töplami $x\to-x$ ga nisbatan simmetrik ($-\frac\pi8-\frac{\pi k}4=\frac\pi8+\frac{\pi(-k-1)}4$), $[-\pi;\pi]$ ham simmetrik, $0$ ildiz emas — yiğindi $0$.
\Javob{a) $\frac\pi8$; \quad b) 0}

%================= 38
\newpage
\Savol{38}
Maksimum $(2;3)$ da: $x=-b=2\Rightarrow b=-2$, \ $c=3$, \ $a<0$. \ $f(6)=4a+3=1\Rightarrow a=-\frac12$.\\
$x<2$ da $f(x)=-\frac12(2-x)+3=\frac x2+2$. Parallellik: $k=\frac12$; \ $g(6)=3+l=1\Rightarrow l=-2$.

a) \[3\cdot\frac{a+b+c}{k+l}=3\cdot\frac{-\frac12-2+3}{\frac12-2}=3\cdot\frac{\frac12}{-\frac32}=-1.\]
b) $y=\frac x2+2$ va $y=\frac x2-2$ parallel töğri çiziqlar orasidagi masofa:
\[\rho=\frac{|2-(-2)|}{\sqrt{1+\frac14}}=\frac4{\frac{\sqrt5}2}=\frac{8\sqrt5}5.\]
\Javob{a) $-1$; \quad b) $\frac{8\sqrt5}5$}

%================= 39
\newpage
\Savol{39}
\[f'(x)=\frac{x^4}8+\frac{x^3}2-\frac{x^2}2-2x=\frac x8\left(x^3+4x^2-4x-16\right)=\frac x8(x+4)(x+2)(x-2).\]
Işoralar: $(-\infty;-4)$: $+$, \ $(-4;-2)$: $-$, \ $(-2;0)$: $+$, \ $(0;2)$: $-$, \ $(2;+\infty)$: $+$.

a) Lokal maksimum ($+\to-$): $x=-4$ va $x=0$ — 2 ta.

b) $[-5;2]$ da nomzodlar: $f(-4)$, $f(0)$ va çetlar.
\[f(-4)=-\frac{128}5+32+\frac{32}3-16+2=\frac{46}{15},\quad f(0)=2,\quad f(-5)=-\frac{13}6,\quad f(2)=-\frac8{15}.\]
Eng kattasi $\dfrac{46}{15}$.
\Javob{a) 2; \quad b) $\frac{46}{15}$}

%================= 40
\newpage
\Savol{40}
$f(x)=a(x+6)(x-2)$, \ $f(-4)=-12a=3\Rightarrow a=-\frac14$. \ $g$: $(-4;3)$ va $(2;0)$ dan: $k=-\frac12$, $l=1$.

a) $a+b+c=f(1)=-\frac14\cdot7\cdot(-1)=\frac74$, \ $k+l=g(1)=\frac12$: \ $\dfrac{7/4}{1/2}=\dfrac72$.

b) $f(x)-g(x)$ — boş koeffitsiyenti $-\frac14$, ildizlari $-4$ va $2$: \ $f-g=-\frac14(x+4)(x-2)$.
\[S=\frac14\int_{-4}^2(x+4)(2-x)\,dx=\frac14\cdot\frac{6^3}6=9.\]
\Javob{a) $\frac72$; \quad b) 9}

%================= 41
\newpage
\Savol{41}
$\angle C=90^\circ$, $CH=6$ — balandlik, $CL=8$ — bissektrisa. $\triangle CHL$: $\cos\angle HCL=\dfrac{CH}{CL}=\dfrac34$.\\
$\angle ACL=45^\circ$, \ $\angle ACH=90^\circ-\angle A=\angle B$ $\Rightarrow$ $\angle HCL=|45^\circ-\angle B|$, \ $2\angle B=90^\circ\pm2\angle HCL$.
\[\sin2B=\cos2\angle HCL=2\cdot\frac9{16}-1=\frac18.\]
a) $AC=\dfrac{CH}{\sin A}=\dfrac6{\cos B}$, \ $BC=\dfrac6{\sin B}$:
\[S=\frac12\cdot\frac{36}{\sin B\cos B}=\frac{36}{\sin2B}=288.\]
b) $c=AB=\dfrac{2S}{CH}=\dfrac{576}6=96$.
\Javob{a) 288; \quad b) 96}

%================= 42
\newpage
\Savol{42}
\fig{\begin{tikzpicture}[scale=0.85]
\coordinate (A) at (0,0); \coordinate (B) at (2.366,2.366); \coordinate (C) at (4.732,0); \coordinate (D) at (2.366,-2.366);
\coordinate (E) at (1,0); \coordinate (F) at (3.366,-1.366);
\draw[thick] (A)--(B)--(C)--(D)--cycle; \draw[thick] (A)--(C);
\draw[thick] (B)--(E)--(F);
\draw[axs] (E)--(D); \draw[ax] (B)--(F);
\draw pic[draw,angle radius=5pt]{right angle=B--E--F};
\node[left] at (A){\small$A$}; \node[above] at (B){\small$B$}; \node[right] at (C){\small$C$}; \node[below] at (D){\small$D$};
\node[below left] at (E){\small$E$}; \node[below right] at (F){\small$F$};
\end{tikzpicture}}
b) $\triangle ABE$: $\angle BAE=45^\circ$. Sinuslar teoremasi:
\[\sin\angle ABE=\frac{AE\sin45^\circ}{BE}=\frac{\frac{\sqrt2}2}{\sqrt3+1}=\frac{\sqrt2(\sqrt3-1)}4=\frac{\sqrt6-\sqrt2}4=\sin15^\circ\ \Rightarrow\ \angle ABE=15^\circ.\]
a) $\angle BEF=\angle BCF=90^\circ$ $\Rightarrow$ $B,E,F,C$ — $BF$ diametrli aylanada: $\angle EFB=\angle ECB=45^\circ$, demak $EF=EB$.\\
$AC$ — simmetriya öqi: $ED=EB=\sqrt3+1$ (\textcolor{aux}{qizil}), \ $\angle ADE=\angle ABE=15^\circ\Rightarrow\angle EDF=75^\circ$.\\
$\triangle EDF$ — teng yonli ($ED=EF$), $DF=2\,ED\cos75^\circ$:
\[DF=2(\sqrt3+1)\cdot\frac{\sqrt6-\sqrt2}4=\frac{(\sqrt3+1)(\sqrt6-\sqrt2)}2=\frac{3\sqrt2-\sqrt2}2=\sqrt2.\]
\Javob{a) $\sqrt2$; \quad b) $15^\circ$}

%================= 43
\newpage
\Savol{43}
Sakkizburçak tomoni $s$, içki burçagi $135^\circ$: $AB$ tomon $Ox$ bilan $45^\circ$ hosil qiladi, $AH$ va $DE$ vertikal.\\
$A$ ning ordinatasi $\dfrac s{\sqrt2}$, abssissasi $u$; \ $E$: abssissasi $u+\dfrac s{\sqrt2}+s+\dfrac s{\sqrt2}=u+s(1+\sqrt2)$, ordinatasi $\dfrac s{\sqrt2}+s$.\\
$A$ grafikda: $\dfrac s{\sqrt2}=\sqrt u\Rightarrow u=\dfrac{s^2}2$. \ $E$ grafikda:
\[s^2\left(1+\frac1{\sqrt2}\right)^2=\frac{s^2}2+s(1+\sqrt2)\ \Rightarrow\ s^2\left(\frac32+\sqrt2\right)-\frac{s^2}2=s(1+\sqrt2)\ \Rightarrow\ s^2(1+\sqrt2)=s(1+\sqrt2).\]
a) $s=1$.

b) $AE$: gorizontal farq $1+\sqrt2$, vertikal farq $s=1$:
\[AE=\sqrt{(1+\sqrt2)^2+1}=\sqrt{4+2\sqrt2}.\]
\Javob{a) 1; \quad b) $\sqrt{4+2\sqrt2}$}

%================= 44
\newpage
\Savol{44}
\fig{\begin{tikzpicture}[scale=0.5]
\coordinate (P) at (-6,0); \coordinate (Q) at (6,0); \coordinate (S) at (0,6);
\draw[thick] (P)--(S)--(Q)--cycle;
\draw[thick,draw=okgreen] (-4,0) rectangle (4,2);
\draw[ax] (S)--(0,0);
\node[above] at (S){\small$S$};
\node[aux,below] at (2,0){\scriptsize$r$}; \node[aux,right] at (4,1){\scriptsize$h$}; \node[aux,right] at (0,4){\scriptsize$6-h$};
\node[below] at (0,-0.6){\small Kesim (asos tomoniga parallel, balandlik orqali)};
\end{tikzpicture}}
Asos tomoni $a=\sqrt{144}=12$, \ $H=\dfrac{3\cdot288}{144}=6$.\\
Silindr balandligi $h$, radiusi $r$: ustki aylana $h$ balandlikdagi kvadrat kesimga içki çizilgan. Öxşaşlikdan $\dfrac r6=\dfrac{6-h}6\Rightarrow r=6-h$.

a) $V(h)=\pi(6-h)^2h$, \ $V'=\pi(6-h)(6-3h)=0\Rightarrow h=2$, $r=4$: \ $V_{\max}=3\cdot16\cdot2=96$.

b) $S(h)=2\pi(6-h)h$ — parabola, maksimum $h=3$ da: \ $S_{\max}=2\cdot3\cdot3\cdot3=54$.
\Javob{a) 96; \quad b) 54}

%================= 45
\newpage
\Savol{45}
Asosiy qarz har oy teng: $\dfrac{60\,000\,000}{24}=2\,500\,000$. Oylik foiz $\dfrac{20\%}{12}=\dfrac1{60}$, qolgan qarzga hisoblanadi.

a) 10-oy boşida qolgan qarz: $60\,000\,000-9\cdot2\,500\,000=37\,500\,000$, foiz $\dfrac{37\,500\,000}{60}=625\,000$:
\[2\,500\,000+625\,000=3\,125\,000.\]
b) Qolgan qarzlar: $2\,500\,000\cdot(24+23+\dots+1)=2\,500\,000\cdot300$. Jami foiz:
\[\frac{2\,500\,000\cdot300}{60}=12\,500\,000.\]
\Javob{a) 3\,125\,000; \quad b) 12\,500\,000}

\ShowJavoblar
\end{document}
